% GENERATED FILE. Edit canonical lessons, then run npm run generate:books.
% canonical-source-hash: fnv1a64:305d796f6047c712
% canonical-lessons: TE-C294-salaha, TE-C294-paiga, TE-C294-annattu, TE-C294-dadapu, TE-C294-mottaniki

\chapter{Advice, Moreover, By the way, Almost, All in all}
\label{ch:te-advice-moreover-by-the-way-almost-all-in-all}
\hlchaptermodality{294}

\begin{chapteropening}
\textbf{By the end of this chapter:} \emph{I can say advice, moreover, by the way, almost and all in all.}

Complete the last lesson of chapter 294: I can say advice, moreover, by the way, almost and all in all.
\end{chapteropening}

\section[\texorpdfstring{\te{సలహా} (salahā)}{సలహా (salahā)}]{\te{సలహా} (salahā) — advice}

\label{lesson:TE-C294-salaha}



\begin{quote}
Before the new one: say the Telugu for the village council, then the Telugu for a loft.
\end{quote}

\subsection*{You'll want to know: \te{సలహా}}
\textbf{\te{సలహా}} — \emph{salahā} — \textquotedblleft{}advice\textquotedblright{}.

\begin{grammarlens}[title={What you've built}]
One more word for this chapter.
\end{grammarlens}

\subsection*{Your turn}
\begin{itemize}
\raggedright
  \item \emph{Say it:} \emph{salahā}
  \item \emph{Say it:} \emph{salahā}, once more
  \item \emph{Say it:} say \emph{salahā}
  \item \emph{Recall:} say the Telugu for the village council, then the Telugu for a loft, then say \emph{salahā} again
\end{itemize}

\begin{tcolorbox}[breakable,colback=teal!4,colframe=teal!35!black,title={Before you move on}]
What does \te{సలహా} mean? (\textquotedblleft{}Advice\textquotedblright{}.) Say it once more.
\end{tcolorbox}
\section[\texorpdfstring{\te{పైగా} (paigā)}{పైగా (paigā)}]{\te{పైగా} (paigā) — moreover, besides}

\label{lesson:TE-C294-paiga}



\begin{quote}
Before the new one: say the Telugu for a loft, then the Telugu for advice.
\end{quote}

\subsection*{You'll want to know: \te{పైగా}}
\textbf{\te{పైగా}} — \emph{paigā} — \textquotedblleft{}moreover, besides\textquotedblright{}.

\begin{grammarlens}[title={What you've built}]
One more word for this chapter.
\end{grammarlens}

\subsection*{Your turn}
\begin{itemize}
\raggedright
  \item \emph{Say it:} \emph{paigā}
  \item \emph{Say it:} \emph{paigā}, once more
  \item \emph{Say it:} say \emph{paigā}
  \item \emph{Recall:} say the Telugu for a loft, then the Telugu for advice, then say \emph{paigā} again
\end{itemize}

\begin{tcolorbox}[breakable,colback=teal!4,colframe=teal!35!black,title={Before you move on}]
What does \te{పైగా} mean? (\textquotedblleft{}Moreover, besides\textquotedblright{}.) Say it once more.
\end{tcolorbox}
\section[\texorpdfstring{\te{అన్నట్టు} (annaṭṭu)}{అన్నట్టు (annaṭṭu)}]{\te{అన్నట్టు} (annaṭṭu) — by the way}

\label{lesson:TE-C294-annattu}



\begin{quote}
Before the new one: say the Telugu for advice, then the Telugu for moreover.
\end{quote}

\subsection*{You'll want to know: \te{అన్నట్టు}}
\textbf{\te{అన్నట్టు}} — \emph{annaṭṭu} — \textquotedblleft{}by the way\textquotedblright{}.

\begin{grammarlens}[title={What you've built}]
One more word for this chapter.
\end{grammarlens}

\subsection*{Your turn}
\begin{itemize}
\raggedright
  \item \emph{Say it:} \emph{annaṭṭu}
  \item \emph{Say it:} \emph{annaṭṭu}, once more
  \item \emph{Say it:} say \emph{annaṭṭu}
  \item \emph{Recall:} say the Telugu for advice, then the Telugu for moreover, then say \emph{annaṭṭu} again
\end{itemize}

\begin{tcolorbox}[breakable,colback=teal!4,colframe=teal!35!black,title={Before you move on}]
What does \te{అన్నట్టు} mean? (\textquotedblleft{}By the way\textquotedblright{}.) Say it once more.
\end{tcolorbox}
\section[\texorpdfstring{\te{దాదాపు} (dādāpu)}{దాదాపు (dādāpu)}]{\te{దాదాపు} (dādāpu) — almost, nearly}

\label{lesson:TE-C294-dadapu}



\begin{quote}
Before the new one: say the Telugu for moreover, then the Telugu for by the way.
\end{quote}

\subsection*{You'll want to know: \te{దాదాపు}}
\textbf{\te{దాదాపు}} — \emph{dādāpu} — \textquotedblleft{}almost, nearly\textquotedblright{}.

\begin{grammarlens}[title={What you've built}]
One more word for this chapter.
\end{grammarlens}

\subsection*{Your turn}
\begin{itemize}
\raggedright
  \item \emph{Say it:} \emph{dādāpu}
  \item \emph{Say it:} \emph{dādāpu}, once more
  \item \emph{Say it:} say \emph{dādāpu}
  \item \emph{Recall:} say the Telugu for moreover, then the Telugu for by the way, then say \emph{dādāpu} again
\end{itemize}

\begin{tcolorbox}[breakable,colback=teal!4,colframe=teal!35!black,title={Before you move on}]
What does \te{దాదాపు} mean? (\textquotedblleft{}Almost, nearly\textquotedblright{}.) Say it once more.
\end{tcolorbox}
\section[\texorpdfstring{\te{మొత్తానికి} (mottāniki)}{మొత్తానికి (mottāniki)}]{\te{మొత్తానికి} (mottāniki) — all in all, in short}

\label{lesson:TE-C294-mottaniki}



\begin{quote}
Before the new one: say the Telugu for by the way, then the Telugu for almost.
\end{quote}

\subsection*{You'll want to know: \te{మొత్తానికి}}
\textbf{\te{మొత్తానికి}} — \emph{mottāniki} — \textquotedblleft{}all in all, in short\textquotedblright{}.

\begin{grammarlens}[title={What you've built}]
One more word for this chapter.
\end{grammarlens}

\subsection*{Your turn}
\begin{itemize}
\raggedright
  \item \emph{Say it:} \emph{mottāniki}
  \item \emph{Say it:} \emph{mottāniki}, once more
  \item \emph{Say it:} say \emph{mottāniki}
  \item \emph{Recall:} say the Telugu for by the way, then the Telugu for almost, then say \emph{mottāniki} again
\end{itemize}

\begin{tcolorbox}[breakable,colback=teal!4,colframe=teal!35!black,title={Before you move on}]
What does \te{మొత్తానికి} mean? (\textquotedblleft{}All in all, in short\textquotedblright{}.) Say it once more.
\end{tcolorbox}
